\section{Experiments}
\label{sec:experiments}

\subsection{Datasets and protocols}

\paragraph{Self-collected (12 participants, 143 trials, 2 526 evaluation
fixations).}
We recruited twelve participants (5 male, 7 female; ages 18--25). All
sessions used a right-eye monocular Pupil Labs Core tracker, a 24-inch
1920$\times$1080 display, and a fixed seating arrangement at 70~cm
viewing distance with the eye tracker's built-in 5-point pre-calibration.
Each trial consists of an 18-point fitting phase followed by a 32-point
testing phase. After standard outlier filtering, 143 trials remain.
Figure~\ref{fig:dataset}(a) shows one example trial.
Figure~\ref{fig:dataset}(b) shows that two participants (Liu Jiaqi and
ZAnna) account for over 90\% of the trials --- we therefore report both
sample-weighted and subject-mean aggregates.
Figure~\ref{fig:dataset}(c) shows that per-trial bias vectors are
direction-stable and of magnitude $25$--$60$~px,
motivating online recalibration. Figure~\ref{fig:dataset}(d)
illustrates the residual vector field after classical calibration on a
single trial.

\paragraph{JuDo1000 (150 participants, target-disjoint split).}
We use the public JuDo1000 dataset~\citep{makowski2020judo} with the
target-disjoint split released by~\citet{anon2026mbnr}: 33 training
target points are used to fit \emph{all} classical baselines, 7 are
held out for validation, and 8 for test. This protocol guards against
the trivial interpolation that would otherwise rank the smoothest RBF
as ``best on training points'' simply because it interpolates them
exactly.

\begin{figure}[t]
\centering
\includegraphics[width=\linewidth]{fig03_dataset.pdf}
\caption{Self-collected dataset structure. (a) One representative trial
of 32 test fixations on a 1920$\times$1080 screen; arrows go from raw
gaze means to targets. (b) Per-subject distribution of classical-baseline
error (sim+RBF), 12 subjects sorted by median; bold orange box outlines
mark the two subjects (Liu Jiaqi, ZAnna) that contribute over 90\% of
trials. (c) Per-trial bias vectors after classical calibration: most
trials retain a coherent bias of 25--60~px that is direction-stable per
session. (d) Residual vector field for a single trial; the spatial
structure is small enough that an additive correction estimated from a
few extra fixations is sufficient.}
\label{fig:dataset}
\end{figure}

\paragraph{Leave-one-subject-out (LOSO).}
For the self-collected dataset, we evaluate every method by
leave-one-subject-out across the 12 participants. For each held-out
subject, the remaining 11 subjects' rows form a training pool from
which 10\% of trials are randomly held out for validation.
The shrinkage MLP is trained from scratch on each fold and evaluated on
the held-out subject without any further adaptation.

\subsection{Honest neural refinement does not beat classical baselines}

Table~\ref{tab:judo} compares all variants on JuDo1000. The strongest
classical baseline is the similarity transform at $21.66$~px; the raw
eye tracker is $21.96$~px. The honest paper method (top-$M$ selection
plus Gaussian noise injection plus the $[64,32,16]$ BN MLP) reaches
$30.46$~px --- $9$~px \emph{worse} than the similarity baseline. An
oracle that ranks baselines on the test set itself reaches $31.94$~px.
The smaller, anchored ``v2'' MLP (Section~\ref{sec:audit}) ties the
similarity baseline at $21.70$~px without surpassing it. The leaky
pipeline reproduces the published $5.82$~px (best-epoch) headline,
confirming that the leak fully accounts for the original gain.

\begin{table}[t]
\caption{JuDo1000 (point-disjoint split, 1\,927 test samples). All
honest variants of multi-baseline neural refinement match or
under-perform the similarity baseline. The leaky pipeline reproduces
the published $5.82$~px.}
\label{tab:judo}
\centering
\small
\begin{tabular}{lrrrr}
\toprule
Method & Mean L2 (px) & Median & Std & p95 \\
\midrule
\multicolumn{5}{l}{\emph{Classical baselines (per-trial fit, point-disjoint)}}\\
origin\_gaze (raw)              & 21.96 & 18.08 & 15.96 & 54.32 \\
similarity                      & \textbf{21.66} & 17.81 & 15.68 & 52.77 \\
poly                            & 25.23 & 21.02 & 17.63 & 59.90 \\
sim+RBF (s=2.0)                 & 47.91 & 44.76 & 23.81 & 89.67 \\
TPS                             & 76.43 & 62.83 & 45.82 & 161.56 \\
\midrule
\multicolumn{5}{l}{\emph{Honest neural refinement (this paper)}}\\
top-$M{=}4$, $\sigma{=}20$, [64,32,16] BN     & 30.46 & 27.5  & 17.6 & 62.7 \\
top-$M{=}4$, $\sigma{=}0$,  [64,32,16] BN     & 31.77 & 29.9  & 15.9 & 62.1 \\
all 11 baselines, $\sigma{=}20$ & 44.39 & 31.3  & 52.4 & 111.9 \\
oracle top-$M{=}4$ on test risk & 31.94 & 30.0  & 15.9 & 62.2 \\
v2 anchored residual MLP        & 21.70 & 17.8  &   -- & 51.9 \\
\midrule
\multicolumn{5}{l}{\emph{Leaky pipeline (matches published; for diagnosis only)}}\\
top-$M{=}4$, leaky inputs       & \textcolor{red}{4.21} & 3.6 & 2.3 & 9.1 \\
\bottomrule
\end{tabular}
\end{table}

Table~\ref{tab:self_collected} shows the same picture on the
self-collected dataset. The strongest classical baseline (sim+RBF) is
$45.53$~px; six honest neural variants land within $\pm 1.0$~px. The
leaky pipeline produces a training-side L2 near $12$~px and the
$42.3$~px test number reported in the original paper.

\begin{table}[t]
\caption{Self-collected dataset (12 subjects, 506 row-random test rows).
Six honest neural variants fall within $\pm 1.0$~px of the strongest
classical baseline (sim+RBF). The published $42.3$~px headline is
reproducible only with the leaky pipeline.}
\label{tab:self_collected}
\centering
\small
\begin{tabular}{lrr}
\toprule
Method & Test L2 (px) & $\Delta$ vs sim+RBF \\
\midrule
origin\_gaze (raw)              & 67.25 & +21.72 \\
similarity                      & 50.45 & +4.92  \\
poly                            & 48.88 & +3.35  \\
GPR                             & 48.29 & +2.76  \\
TPS                             & 49.58 & +4.05  \\
sim+RBF (anchor)                & \textbf{45.53} & 0.00 \\
\midrule
uniform mean of 6 baselines     & 45.78 & +0.25  \\
per-trial LOO BMA, $T{=}10$     & \textbf{44.62} & --0.91 \\
v2 anchor-residual MLP          & 44.71 & --0.82 \\
v2 stacker, LOSO sample-weighted& 45.03 & +0.34 \\
DeepSets context refiner, $K{=}12$ LOSO & 48.84 & +3.31 \\
\midrule
honest top-$M{=}4$, [64,32,16] BN, $\sigma{=}20$ (paper-faithful) & 45.06 & --0.47 \\
\midrule
leaky pipeline (matches published) & \textcolor{red}{42.3} & --3.2 \\
\bottomrule
\end{tabular}
\end{table}

\subsection{Online shrinkage is the win}

Figure~\ref{fig:shrinkage}(a) plots test L2 against context size $K$
for three policies. Raw mean-bias correction (orange) hurts at $K=1$
and only crosses below the anchor at $K \ge 5$. The globally tuned
fixed-$\lambda$ policy (magenta) improves over raw mean by removing the
small-$K$ noise penalty. The learned per-trial shrinkage (green)
matches fixed-$\lambda$ for very small $K$ and substantially exceeds it
once $K \ge 12$, achieving $42.29$~px (a $17\%$ reduction over the
$44.69$~px anchor) at $K=12$ and $40.20$~px ($10\%$ reduction) at
$K=18$.

\begin{table}[t]
\caption{Test mean L2 (px) on the self-collected dataset across context
sizes $K$. At small $K$, raw mean-bias correction (column ``raw $\lambda{=}1$'')
is dominated by sampling noise. The fixed-$\lambda$ sweep recovers most
of the gain. The learned shrinkage MLP (LOSO) matches the fixed sweep
at small $K$ and dominates it at $K \ge 12$.}
\label{tab:shrinkage}
\centering
\small
\begin{tabular}{rrrrr}
\toprule
$K$ & raw $\lambda{=}1$ & fixed $\lambda^\star$ ($\lambda^\star$) & learned $\lambda$ (LOSO) & $\Delta$ vs anchor (44.69)\\
\midrule
0  & 44.69 & 44.69 & 44.69 & 0.00 \\
1  & 58.80 & 43.26 (0.2) & 63.88 & --1.43 (fixed) \\
2  & 48.98 & 42.74 (0.4) & 60.08 & --1.95 (fixed) \\
3  & 45.99 & 41.97 (0.4) & 61.29 & --2.72 (fixed) \\
5  & 43.27 & 41.02 (0.4) & 56.32 & --3.67 (fixed) \\
8  & 40.61 & 39.37 (0.6) & 54.13 & --5.32 (fixed) \\
12 & 37.95 & 37.52 (0.8) & \textbf{42.29} & \textbf{--7.17 (fixed)}, \textbf{--2.40 (LOSO)} \\
18 & 38.60 & 38.18 (0.8) & \textbf{40.20} & --6.51 (fixed), \textbf{--4.49 (LOSO)} \\
\bottomrule
\end{tabular}
\end{table}

The learned shrinkage policy underperforms fixed-$\lambda$ for
$K \le 8$. This is consistent with the small training-trial count for
small-$K$ episodes in our LOSO folds; with more training subjects we
expect the gap to close. We discuss this in Section~\ref{sec:discussion}.

\subsection{Per-subject LOSO}

Figure~\ref{fig:loso} reports per-subject results at $K{=}12$. Three
subjects --- the only three with enough trials for $K{=}12$ episodes ---
are evaluated and all three improve under the learned shrinkage:
\textbf{Liu Jiaqi} ($40.79 \to 35.53$~px, $-13\%$),
\textbf{ZAnna} ($45.90 \to 44.02$~px, $-4\%$),
\textbf{wzx} ($57.16 \to 47.33$~px, $-17\%$).
The remaining nine subjects have insufficient trials to form $K{=}12$
context-plus-query episodes (see Discussion). Sample-weighted mean
across the three viable folds: $-5.7$~px ($-12\%$).

\begin{figure}[t]
\centering
\includegraphics[width=\linewidth]{fig07_loso.pdf}
\caption{Per-subject leave-one-subject-out evaluation at $K=12$
context fixations. The shrinkage policy is trained on the other 11
subjects and applied to the held-out subject. All three subjects with
enough trials to admit $K=12$ episodes (Liu Jiaqi, ZAnna, wzx) improve;
the remaining nine subjects are flagged ``insufficient ctx''.}
\label{fig:loso}
\end{figure}

\subsection{Ablations}

We performed comprehensive ablations on JuDo1000 along six axes:
selection strategy ($\{$all, top-$M\!\in\!\{1{,}2{,}3{,}4{,}6{,}8\}$,
oracle$\}$); noise $\sigma_{\max}\in\{0,5,10,20,40,80\}$~px; network
capacity ($\{[16],[64],[64,32],[64,32,16],[256,128],[1024,512]\}$);
batch normalization on/off; baseline pool drop-one; and seeds
$\{1047, 7, 42, 123, 2024\}$. Results land within $\pm 1$~px of each
other and never reach the strongest classical baseline at $21.66$~px.
Full numbers in Appendix~\ref{app:ablations}.

For online shrinkage, the central knobs are $K$ (Table~\ref{tab:shrinkage})
and the choice of anchor (sim+RBF dominates similarity, polynomial,
GPR by 1--2~px). Replacing the learned $\lambda$ with a constant
$\lambda \in \{0, 0.5, 1\}$ recovers the fixed-$\lambda$ sweep curve;
replacing the per-trial bias estimate with a global subject-mean bias
loses all benefit, confirming that the per-trial component is essential.
